\documentclass[11pt]{article}
\usepackage[margin=1in]{geometry}             
	\geometry{letterpaper}
	\usepackage{pdfpages}
\usepackage{setspace}
\usepackage{parskip}		%% Adjusts paragraph settings
\usepackage{graphicx}
\usepackage{enumitem}				%% For numbered or bulleted lists
	%\setlist[itemize]{leftmargin=20mm}	%% Change list margin



\begin{document}

\textbf{Boos and Pascale, 2021 }
\begin{itemize}
    \item sensible heat fluxes from orography into upper layers of atmosphere draw water vapor upslope of Sierra Madres which then condense and precipitate. Near-surface air flows upslope during daytime and downslope at nighttime.
    \item Understanding of blocking of easterlies by Sierra Madre Occidentals on NAM circulation has been limited by low-resolution GCMs that do not fully resolve the SMOs (or perhaps the SMOs are too low, which we would be testing).
    \item time-mean vertically-integrated moisture flux convergence in core NAM is driven by time-mean winds (not transients, i.e. diurnal cycle)
    \item used a 0.25\textdegree{} model, SMO are 3km in control
    \item In flat mexico simulation, at 700 hPa, easterly winds (flowing westward) exist over Mexico. In control, 700 hPa winds are westerly (flowing eastward). This is confirmed with the stationary wave model. So, the stationary wave induced by the orography is what drives the time-mean westerly winds that drive upslope flow up the SMOs.
    \item isenotropes tilt downward to the north in southwest north america due to temperature maxima $\sim$38\textdegree{N}. This leads adiabatic zonal flow to turn southward when it encounters the SMO, because if it turned northward it would run into the ground? This generates anticyclonic circulation southwest of Mexico. This anticyclone is weak in low-res models. 
    \item warm and moist (high MSE) air mass exists over the GoC. Due to diurnal cycle of solar heating over land, in afternoon there is also a peak of MSE over western slopes of SMO. So, when GoC airmass flows eastward due to mechanically-forced stationary wave, its MSE is further increased by daytime sensible heat fluxes over the mountain while also cooling (is cooling ``release of CAPE") due to upslope flow. This generates moist convection. So, maybe we don't see that much of a stronger NAM in E3 because of poor resolution of the GoC and that initially high MSE airmass, even despite increased eastward flow or something.
    \item Analysis of seasonal cycle shows that upslope (eastward) flow peaks in spring, but MSE is not high enough then so moist convection/precip doesn't occur. Peak precipitation occurs when MSE peaks and upslope flow is still relatively strong. \textbf{So, ultimately, while thermodynamic cycle is important in driving seasonal increase in MSE (and CAPE), mechanical impact of SMO on stationary wave is necessary to turn it into precipitation.}
    \item Based on their response to Kau and Marcus, it seems how well the model represents the deflection of the westerly jet to the north and/or south of the Rockies will determine how much the zonal mean flow interacts with the Sierra Madre Occidental to produce the stationary wave. I'm not sure how well FLOR vs. E3 represents this, and it will certainly be impacted by the MAX topo sims that also elevate the Rockies. 
    \item thermal forcing (rather than mechanical forcing) might impact transients that are relevant for moisture convergence north of the NAM core (i.e. the southwestern U.S.)
\end{itemize}
\newpage

\textbf{Bhattacharya et al., 2023 ERL}
\begin{itemize}
    \item They define/quantify the \textit{observed} relationship between CA margin (25\textdegree{N}--35\textdegree{N}, 125\textdegree{W}--110\textdegree{W}) SST (\textbf{\textit{monthly}} anomalies from ICOADS 1\textdegree{} product -- Ishii et al., 2005) and extreme precipitation (95th percentile of \textbf{\textit{daily}} precipitation rates in GPCC 1\textdegree{} product). Used same data to calculate probability of summer with $>$6 days of extreme precipitation ($>$p95) when CA margin temperatures were anomalously warm ($>$1$\sigma$).
    \item Tested quantifying the relationship between CA margin SST and monthly mean precipitation from multiple products -- 0.25\textdegree{} GPCC, CPC, GPCP, and NARR reanalysis). Did this focus on monthly mean precipitation since monthly mean data is what was available from CMIP5 and CMIP6 output
    \item Then, evaluate ability of CMIP5 and CMIP6 models to reproduce the observed CA margin - extreme northern NAM (28\textdegree{N}--37\textdegree{N}, 115\textdegree{W}--108\textdegree{W}) precip relationship. Analyzed historical model simulations. Show composite output for models that capture the correlation well vs. poorly.
    \item The strongest NAM years occurred during both positive and negative ENSO phases, adding to the body of literature showing inconclusive relationship between tropical Pacific climate and the NAM
    \item CA margin SSTs are positively correlated with southerly/southwesterly moisture flux over Baja and coastal California. This happens (they say) because the warmer SST weaken the southeast edge of the NPSH by reducing the local static stability.
    \item Then, they say, moisture flux into the desert increases humidity and dewpoints, which fuels instability and, ultimately, enhances CAPE. 
    \item \textbf{\textit{Caveat with Clemesha reference, say that though they cannot say with certainty whether outflow from the NAM amplifies CA margin SSTs (through low cloud feedbacks), they argue that the CA margin SST warming occurs first and drives the cyclonic circulation anomaly and moisture flux into the southwest.}}
    \item Interestingly, they find that models with warm SST biases along the CA margin (\textbf{\textit{which maybe exist due to model biases in low cloud simulation, my interpretation not theirs}} perform worst in capturing the correlation between CA margin SSTs and NAM precipitation.
    \item They actually only looked at 2 HighResMIP models, but found that increasing horizontal resolution improved the CA margin SST bias and the strength of the correlation between CA margin SSTs and NAM precipitation. Argue that HR models don't simulate NAM climatology all that well anyways, but they only looked at 2 models and their figure ``demonstrating" this (S9) is trash. They suggest that biases in the simulation of large-scale climate play an important role in models' ability to capture coupling between SST-atm circulation-precip. But, we don't see significant changes in the CA SST in FLOR HiTOPO?
    \item They ultimately argue that the reason models with stronger CA margin SST biases simulate the correlation with NAM precipitation more poorly is because additional warming represents less of a perturbation and therefore is less efficient at altering atmospheric circulation. But, then, do these models have a stronger default NAM? I don't totally understand...
    \item Ultimately, they also find that the relationship is only relevant/strongest for precipitation in the \textit{northern} NAM domain, and is weak for the core region in the Mexico. \textbf{\textit{This is good and possibly conforms with our HighResMIP and FLOR analyses whereby large-scale circulation does not affect NAM biases, rather, direct orographic blocking processes are more important (i.e. even though we improve precip biases in HI-MEX, circulation biases are worsened, thereby circulation must not be important for representing precip). But, as in HI-GBL, reducing NAM biases in the SW US is consistent with reductions in circulation patters.}}
    \item They try to say that, for models that exhibit a strong SST-precip correlation over the historical period, there is a persistence of this strong correlation in the future (as a counter to their hypothesis that as CA margin SSTs warm, the correlation might become weaker given that models with warm SST bias have a weak correlation). But, okay, I really do not understand their analysis of the SST-Precip relationship for future scenarios (figs S12 and S13). To me, it seems that the historical correlation is actually not statistically significant for half of the time period....
\end{itemize}

\subsection{Influence of topography on seasonal feedbacks}
\subsubsection{Soil Moisture}
Low spring soil moisture in SWNA related to stronger subtropical ridge (Vargas Zeppetello et al., 2024), which could increase NAM strength.

Low spring soil moisture in PNW hypothesized to increase subtropical ridging and increase NAM strength -- Hu and Feng (2007) hypothesis. Mike Pye (UIUC) tested this and presented results to the contrary at AMS 2025. 

Higher topography will impact winter and spring precipitation that alters the antecedant soil moisture conditions relevant for these feedbacks.

Ways to possibly explain why FLOR has better NAM and E3 doesn't, perhaps FLOR has a better representation of the seasonal cycle of MSE and/or diurnal cycle of MSE (c.f. Boos and Pascale, 2021 NG). 

Monsoon retreat bias (i.e. it persists too long in October in GCMs) is due to Pacific SST biases, not spatial resolution, due to underestimation of North Pacific Subtropical High strength (Ye and Wang, 2023 \textit{J. Clim.})

\subsubsection{Rocky Mountain snowpack}
Some studies have argued that there is a reduction in NAM strength during years with a deep/extensive Rocky Mountain snowpack due to albedo feedbacks on the subtropical ridge strength. Essentially, during years of deeper snowpack, there is stronger albedo cooling of the surface and the subtropical ridge does not extend as far into North America. This suppresses the monsoon in SWNA. Higher topography should have an impact on this.

Though, coarse resolutions may underrepresent future changes in snowpack and thus this feedback (Hsieh et al., 2025).

Related references:
\begin{itemize}
    \item Notaro and Zarrin, 2011 \textit{GRL} https://doi.org/10.1029/2011GL048803
    \item Hsieh et al., 2025 \textit{Climate Dynamics} https://doi.org/10.1007/s00382-025-07651-6
    \item Cleveland et al., 2025 \textit{J. Geophysical Research Atmospheres} https://doi.org/10.1029/2024JD043026
\end{itemize}

this will be a helpful reference for understanding increased moisture in eastern rocky mountains: https://rmets.onlinelibrary.wiley.com/doi/full/10.1002/joc.8121


\end{document}